\section{Where the number lives}
\label{sec:vslives}

We can compute the fair strike three ways and get one answer per
model.  \Cref{fig:vspins} showed that different models give different
answers.  This section locates the disagreement on the strike axis,
and the location explains everything.

The tool is the accrual function.  Write $\mathcal{S}(k)$ for the
share of the strike-side integral \eqref{eq:vsrepl} accrued by strikes
up to $k$: run the integral from the far left to $k$, divide by the
whole.  Then $\mathcal{S}$ climbs from $0$ to $1$, and reading it at
two strikes tells you what fraction of the fair strike lives between
them.

\begin{figure}[htbp]
\centering
\paperfig{\textwidth}{fig_vs_share.pdf}
\caption{Where the fair strike accrues.  (a)~The accrual share
$\mathcal{S}(k)$ on the running slice, with the quoted span shaded and
its endpoints dotted: the span carries \MacVsShareQuotedPct\% of the
integral, so \MacVsShareBeyondPct\% accrues where there are no quotes
at all.  (b)~The beyond-quotes share for all sixteen frozen nodes
(the snapshot's stored fits), against maturity: from
\MacVsShareGalleryMinPct\% to \MacVsShareGalleryMaxPct\%, growing with
maturity for SPY.}
\label{fig:vsshare}
\end{figure}

\Cref{fig:vsshare}(a) draws $\mathcal{S}$ on the running slice, with
the quoted span shaded.  Follow the curve from the left.  It is
essentially flat until $k\approx-1$: the very deep puts, for all their
$1/K^{2}$ weighting, carry prices too small to matter.  It then climbs
steeply through the belly --- $|k|\le0.10$ alone accrues
\MacVsShareAtmPct\% of the number --- and flattens again on the call
side, where prices die quickly.  The two dots mark the edges of the
quoted span: between them the market's quotes constrain the integrand
directly, and that stretch carries \MacVsShareQuotedPct\% of the
integral.  The remaining \MacVsShareBeyondPct\% accrues at strikes
where no quote exists --- where the integrand is whatever the fitted
model's continuation says it is.  About a fifth of this traded number
is, on this node, pure extrapolation.

Panel~(b) repeats the measurement across the whole frozen gallery ---
all sixteen (underlier, expiry) nodes of the book's snapshot, each
read from the fit it was published with (the snapshot's stored fits;
\cref{app:vsprotocol} records the provenance) and each with its own
quoted span.  The beyond-quotes share runs from
\MacVsShareGalleryMinPct\% to \MacVsShareGalleryMaxPct\%.  The
pattern is worth reading.  What protects a node is not how many
quotes it has but how many \emph{widths of the smile} its quoted span
covers.  A strike matters through its standardized distance from the
money, $k/(\sigma\sqrt\tau)$ with $\sigma$ that expiry's ATM implied
volatility: the smile's natural width is $\sigma\sqrt\tau$, and a
quoted span of roughly fixed width in $k$ covers fewer of those
widths as $\sqrt\tau$ grows.  SPY's long-dated spans cover fewer, and
the unquoted share climbs steadily with maturity toward a quarter of
the number; the marked running node sits at \MacVsShareBeyondPct\%.
NVDA's quoted spans are generous relative to its volatility and its
share stays under ten percent.

Now put the two exhibits of the chapter together.
\Cref{fig:vspins}(b) showed three models disagreeing on the fair
strike by up to \MacVsPinsMaxGapBp\ vol bp while agreeing on the
quoted span to \MacVsPinsBellyAgreeBp\ bp.  \Cref{fig:vsshare} says
the disagreement has exactly one place to live: the
\MacVsShareBeyondPct\% of the integral that accrues beyond the
quotes.  The same comparison runs across model chapters: Chapter~4's
whole-surface sheet, fitted to the same mid quotes, implies
$\sigma_{\rm vs}=\MacVsFieldSvsPct\%$ at this expiry ---
\MacVsLvGapBp\ vol bp from the running slice's
\MacVsThreeSvsPct\%.  None of these models is wrong.  Each reprices
the quotes; the quotes simply do not decide the number.  A fit error
of a few basis points \emph{per quote} coexists with a fair-strike
spread ten times larger, because the fair strike reads a region the
fit error never visits.

This is the chapter's instance of a distinction the book keeps
returning to.  What the quotes pin --- the belly, the well-quoted
central stretch of strikes --- is information.  What they leave free
--- the wings' share --- is not noise around zero; it is a quantity
the model must supply, and different structural choices supply
different amounts.  \Cref{sec:vsceiling,sec:vsenvelope} will show
that the freedom is bounded: no admissible continuation --- none that
still prices a genuine law --- can push the share arbitrarily high.
But within those bounds, the number is chosen, not estimated.

One consequence deserves a paragraph before we turn to the bounds,
because it changes how a desk should read the instrument.  A traded
var-swap quote is one \emph{observation} of precisely the combination
of unquoted prices the listed strikes leave loose.  Chapter~2's
calibration objective (\cref{sec:calibration}) accommodates it
directly: the model's own fair strike --- computed by whichever route
of \cref{sec:vsthree} is native --- becomes one more residual against
the quoted level, in volatility units, with a weight the modeler
declares; the derivative of $w_{\rm vs}$ in the model's parameters
comes in closed form by the linearity of \cref{sec:vsthree}
(Chapter~2 already recorded it as \cref{eq:varswaprow}).  Weighted
softly, the quote bargains with the options rather than overruling
them, and it spends its influence exactly where \cref{fig:vsshare}
says its information lives: in the wings.  A desk that has such a
quote should use it; the sections that follow are addressed to the
more common case where it does not.

We now know where the number lives and why fitted models disagree
about it.  The next two sections establish what no model, and no
market, can do in that region.
